\subsection{Déformation harmonique pondérée et sélection de la branche superstable}
\label{mab:sec:extension-dodecafasica-ponderada-armonica}

Le profil uniforme possède une extension paramétrique qui conserve les douze
phases et sépare la déformation harmonique du paramètre dynamique. Pour
\(1\le m\le12\), soient

\begin{equation}
 \phi_m=\frac{(30m-10)\pi}{180},
 \qquad
 p(\phi)=12\cos(3\phi)-\cos(4\phi).
 \label{mab:eq:fases-y-deformacion-armonica-armonizada}
\end{equation}

Pour tout \(\eta\) réel tel que les poids restent positifs, nous définissons

\begin{equation}
 w_m(\eta)=\frac{1+\eta p(\phi_m)}{12}.
 \label{mab:eq:pesos-dodecafasicos-armonizados}
\end{equation}

Les sommes des harmoniques de rang trois et quatre sur les douze phases s’annulent ;
par conséquent,

\[
 \sum_{m=1}^{12}w_m(\eta)=1.
\]

Soit

\begin{equation}
 s_\eta^2=\frac{2}{\sum_{m=1}^{12}w_m(\eta)\phi_m^2},
 \qquad
 u_\eta(x)=1-\sum_{m=1}^{12}w_m(\eta)
 \cos(s_\eta\phi_m x).
 \label{mab:eq:familia-ponderada-armonizada}
\end{equation}

Le développement du cosinus et la définition de \(s_\eta\) donnent

\[
 u_\eta(x)=x^2+O(x^4)\qquad(x\to0).
\]

Par conséquent, \(u_\eta\) est réelle--analytique, paire et conserve un point
critique quadratique à l’origine. Nous introduisons le paramètre dynamique par

\begin{equation}
 f_{\lambda,\eta}(x)=1-\lambda u_\eta(x).
 \label{mab:eq:familia-feigenbaum-ponderada-armonizada}
\end{equation}

Un paramètre superstable de période exacte \(2^n\) satisfait

\[
 f_{\lambda,\eta}^{\,2^n}(0)=0,
 \qquad
 f_{\lambda,\eta}^{\,2^k}(0)\ne0\quad(0\le k<n).
\]

L’équation peut posséder plusieurs branches. On fixe une branche
\(\mathscr B\) en parcourant, immédiatement après le paramètre du
niveau précédent, un maillage adaptatif ; on prend le premier intervalle avec changement
de signe et l’on applique une dichotomie. La notation
\(\lambda_n^{\mathscr B}(\eta)\) incorpore ce sélecteur et n’attribue pas
d’unicité globale à l’équation superstable. Le premier niveau est déterminé
exactement :

\[
 \lambda_1^{\mathscr B}(\eta)=\frac1{u_\eta(1)},
\]

si \(u_\eta(1)>0\), parce que
\(f_{\lambda,\eta}^{\,2}(0)=f_{\lambda,\eta}(1)=0\).
Étant donnés trois niveaux consécutifs et

\[
 x_n^{\mathscr B}(\eta)
 :=f_{\lambda_n^{\mathscr B}(\eta),\eta}^{\,2^{n-1}}(0),
\]

on forme

\begin{equation}
 \delta_n^{\mathscr B}(\eta)
 =\frac{\lambda_{n-1}^{\mathscr B}(\eta)-
         \lambda_{n-2}^{\mathscr B}(\eta)}
        {\lambda_n^{\mathscr B}(\eta)-
         \lambda_{n-1}^{\mathscr B}(\eta)},
 \qquad
 \alpha_n^{\mathscr B}(\eta)
 =\frac{|x_{n-1}^{\mathscr B}(\eta)|}
        {|x_n^{\mathscr B}(\eta)|}.
 \label{mab:eq:aproximantes-ponderados-armonizados}
\end{equation}

Pour \(\eta=0\), on retrouve exactement le profil uniforme. L’extension
pondérée introduit un paramètre transversal : elle permet d’étudier
la transversalité de la famille par rapport à la variété stable du
renormalisateur sans transformer une coïncidence finie en une preuve
d’universalité.

\begin{remark}[Contrôle négatif]
L’extension est réfutée si les poids perdent leur positivité dans le domaine
déclaré, si la normalisation cesse de produire un coefficient quadratique
unitaire, si le sélecteur n’exclut pas les périodes diviseurs ou si une branche est choisie
en utilisant comme entrée les valeurs universelles terminales.
\end{remark}
